\documentclass[11pt]{article}
\usepackage[margin=25mm]{geometry}
\usepackage{fontspec}
\setmainfont{Latin Modern Roman}
\usepackage{amsmath,amssymb,mathtools}
\usepackage{unicode-math}
\setmathfont{Latin Modern Math}
\usepackage{xurl,hyperref}
\hypersetup{colorlinks=true,urlcolor=blue}
\setcounter{secnumdepth}{0}
\setlength{\parindent}{0pt}
\setlength{\parskip}{5pt}
\setlength{\emergencystretch}{3em}
\newcommand{\sref}[2]{\hyperlink{source:#1:#2}{[S#2]}}
\newcommand{\cataloguescope}[1]{\paragraph{Recorded target.}#1}
\begin{document}
\section{Decomposition of completions of reloids}\label{statement}

\subsection{Definitions and mathematical statement}

\paragraph{Shared notation.}$\mathbb N=\{1,2,\ldots\}$, and $\mathbb Z,\mathbb Q,\mathbb R,\mathbb C$ denote the integers, rationals, reals and complex numbers. Any other conventions are specified locally.


For a set $X$, let $\mathcal F(X)$ be all filters on $\mathcal P(X)$, including the improper filter $\bot_X=\mathcal P(X)$. A filter is a nonempty family closed under finite intersections and supersets. Order filters by $\mathcal A\sqsubseteq\mathcal B$ when $\mathcal A\supseteq\mathcal B$ as families. Thus $\bigsqcup$ is intersection of filter families and $\bigsqcap$ is the filter generated by their union; the empty join is $\bot_X$ and the empty meet is $\top_X=\{X\}$. For $E\subseteq X$, $\uparrow E=\{U\subseteq X:E\subseteq U\}$ is its principal filter. Write $\mathcal A\mathrel{\#}\mathcal B$ when every $U\in\mathcal A$ and $V\in\mathcal B$ have $U\cap V\ne\varnothing$, equivalently $\mathcal A\sqcap\mathcal B\ne\bot_X$.

A reloid $f:X\to Y$ is a filter on $X\times Y$, ordered by reverse inclusion of filter families. For a relation $R\subseteq X\times Y$, its principal reloid is $\uparrow R$. The reloid $g\circ f:X\to Z$ is generated by $G\circ F$ for $F\in f,G\in g$, where $G\circ F=\{(x,z):\exists y,(x,y)\in F,(y,z)\in G\}$. The reverse $f^{-1}$ is generated by $F^{-1}=\{(y,x):(x,y)\in F\}$ for $F\in f$. The identity $1_X$ is $\uparrow\Delta_X$, where $\Delta_X=\{(x,x):x\in X\}$.For a reloid $f:X\to Y$, define its source-singleton restriction $f|_x$ as the meet of $f$ and $\uparrow(\{x\}\times Y)$. Set
\[
 \operatorname{Compl}(f)=\bigsqcup_{x\in X}f|_x,
 \qquad \operatorname{CoCompl}(f)=\bigl(\operatorname{Compl}(f^{-1})\bigr)^{-1}.
\]
Call $f$ complete if $\operatorname{Compl}(f)=f$, and co-complete if $\operatorname{CoCompl}(f)=f$.
\paragraph{Statement recorded in the supplied source.}
Do all five identities hold for arbitrary reloids whenever their displayed compositions are defined?
\begin{enumerate}
\item $\operatorname{Compl}(g\circ f)=\operatorname{Compl}(g)\circ f$ if $f$ is co-complete.
\item $\operatorname{CoCompl}(f\circ g)=f\circ\operatorname{CoCompl}(g)$ if $f$ is complete.
\item $\operatorname{CoCompl}(\operatorname{Compl}(g)\circ f)=\operatorname{Compl}(g\circ\operatorname{CoCompl}(f))=\operatorname{Compl}(g)\circ\operatorname{CoCompl}(f)$.
\item $\operatorname{Compl}(g\circ\operatorname{Compl}(f))=\operatorname{Compl}(g\circ f)$.
\item $\operatorname{CoCompl}(\operatorname{CoCompl}(g)\circ f)=\operatorname{CoCompl}(g\circ f)$.
\end{enumerate}
In item 2, the typing is $g:Z\to X$ and $f:X\to Y$; in items 1,3,4,5 it is $f:X\to Y$ and $g:Y\to Z$. No common endomorphism assumption is imposed. \sref{statement}{2}\sref{statement}{3}
\subsection{Short English statement}
Determine whether completion and co-completion obey all five stated composition identities, retaining the one-sided completeness hypotheses and correct order of composition.
\subsection{Sources}
\begin{description}
\item[\hypertarget{source:statement:1}{S1}] ULAM.AI and the UnsolvedMath Contributors, UnsolvedMath: A Curated Collection of Open Mathematics Problems, catalogue entry. CC BY 4.0. \url{https://huggingface.co/datasets/ulamai/UnsolvedMath}.
\item[\hypertarget{source:statement:2}{S2}] Open Problem Garden, Decomposition of completions of reloids, catalogue entry. Original problem statement, full mathematical discussion, bibliography and comments. \url{https://www.openproblemgarden.org/op/decomposition_of_completions_of_reloids}.
\item[\hypertarget{source:statement:3}{S3}] Victor Porton, General Topology as Ordered Semigroup Actions, author manuscript, version of 21 July 2026. Completion and co-completion; decomposition conjecture containing the five formulas. \url{https://github.com/vporton/agt-combined/blob/759e45cced7769cf3a84ce60ed6c30f4ce9ffa8e/chap-reloids.tex}.
\end{description}
\label{end:statement}
\end{document}
